% ১. প্রোগ্রামের ধারণা
\section{প্রোগ্রামের ধারণা}

\begin{learningbox}
এই টপিক শেষে শিক্ষার্থী program, instruction, source/object/executable program এবং software-এর সম্পর্ক ব্যাখ্যা করতে পারবে।
\end{learningbox}

\begin{definitionbox}{প্রোগ্রাম}
কোনো নির্দিষ্ট কাজ করানোর জন্য কম্পিউটারকে দেওয়া ধারাবাহিক নির্দেশসমষ্টিকে প্রোগ্রাম বলা হয়।
\end{definitionbox}

\begin{center}
\begin{tabular}{|p{0.24\textwidth}|p{0.34\textwidth}|p{0.28\textwidth}|}
\hline
\textbf{ধরন} & \textbf{অর্থ} & \textbf{উদাহরণ} \\
\hline
Source program & programmer যে code লেখে & \texttt{sum.c} \\
\hline
Object program & translated machine-level file & object code \\
\hline
Executable program & সরাসরি run করা যায় & \texttt{sum.exe} \\
\hline
\end{tabular}
\end{center}

\begin{examplebox}{বাস্তব উদাহরণ}
Calculator app-এ user দুইটি সংখ্যা দিলে যোগফল দেখায়। input নেওয়া, যোগ করা এবং output দেখানোর নির্দেশগুলোই program।
\end{examplebox}

\begin{practicebox}
\begin{enumerate}
\item Program কী?
\item Source program ও object program-এর পার্থক্য লেখ।
\item Program ও software-এর সম্পর্ক ব্যাখ্যা কর।
\end{enumerate}
\end{practicebox}
